\section{Step-GUI}

\subsection{数据构建思路}

我们希望构建的不是一个只会执行固定 GUI 指令的狭义策略模型，而是一个同时保留\textbf{通用多模态能力}、具备丰富世界知识、并拥有\textbf{强 agentic 能力}的视觉语言基础模型。为此，Step-GUI 的训练数据分为中训练（mid-train）与冷启动（cold-start）两类阶段。

\paragraph{中训练数据。}
中训练的目标是为模型建立 GUI 智能体所需的基础能力，包括：保留通用视觉与语言理解、掌握细粒度视觉 grounding、理解轨迹格式与动作协议、形成从自然语言指令到可执行动作的初步映射。我们因此混合两大类数据：
\begin{itemize}
\item \textbf{通用多模态与知识数据}：用于保持语言流畅性、世界知识与通用视觉理解。
\item \textbf{Agent 导向数据}：包括动作对齐数据、轨迹数据和跨平台环境数据，使模型熟悉 Android、Windows、macOS、Ubuntu 等系统中的操作模式。
\end{itemize}

\paragraph{冷启动数据。}
在已有基础能力上，冷启动阶段主要用于\textbf{注入缺失知识}并\textbf{校正执行行为}。我们观察到，很多 GUI 失败并非源于缺少操作轨迹，而是源于知识盲区，例如不了解某应用页面结构、服务语义或任务规则。因此，冷启动阶段会根据轨迹执行错误反推出缺失知识，并将其重写成面向模型的问答或视觉问答样本，同时保留一部分轨迹数据用于稳定动作格式与执行习惯。

\subsection{CSRS：校准式步骤奖励系统}

GUI 任务通常是长轨迹、多步骤、稀疏奖励的执行过程，简单依赖模型自我评价会带来严重幻觉与错误归因。CSRS 的核心思想是：\textbf{不直接相信模型生成的步骤级解释，而是用轨迹级结果信号去校准步骤级数据。}

CSRS 主要由以下部分组成：
\begin{itemize}
\item \textbf{轨迹级验证}：通过自动评测脚本或人工结果标注判断整条任务轨迹是否成功。
\item \textbf{步骤级知识提炼}：对成功轨迹中的关键步骤、决策依据和环境模式做结构化抽取，转化成更可靠的训练样本。
\item \textbf{渐进式训练闭环}：模型不断生成新轨迹，系统不断校准、过滤和回流高质量数据，从而形成自进化循环。
\end{itemize}

与传统逐步标注相比，CSRS 通过“轨迹成功性”这一更高置信度的信号，显著降低了步骤标注成本，同时也让模型能从真实执行结果中学习更稳定的 GUI 推理模式。

\subsection{训练机制}

Step-GUI 使用三阶段的渐进式训练策略，把\textbf{新探索}与\textbf{知识过滤}两条数据流结合起来，不断提升模型能力。训练过程中还引入了强化学习式优化，用于提升长程任务中的规划与执行质量。

\paragraph{半在线探索。}
长程 GUI 任务的奖励往往极其稀疏，因此单纯依赖随机探索难以有效学习。本文采用\textbf{Semi-Online}策略：对于原本未能完成任务的 rollout，系统在第二轮执行时向模型注入有限的 ground-truth 提示，使模型能够接触到更高质量的成功轨迹。这些“带回溯的成功经验”会进一步转化为高价值训练样本。

\paragraph{稳定性增强。}
为了让 RL 训练更稳定，文中还引入了若干增强策略，包括动态探索范围、基于重要性采样的样本复用，以及对 clipping 后梯度完全消失问题的修正。其目标是在不过度牺牲稳定性的前提下，保留对低概率但有价值动作的学习能力，避免模型过早塌缩到局部最优。

\subsection{模型定位}

Step-GUI 的目标并不是仅在某一个 GUI 基准上刷分，而是形成一类可扩展的 GUI 专家模型。4B 版本强调\textbf{本地部署与隐私友好}，适合设备侧执行；8B 版本则在 grounding、长程任务完成率和多模态综合能力之间取得更优平衡。实验表明，这种设计既提高了 GUI 专项能力，也没有明显损害模型在 OCR、图表理解、一般视觉问答等任务上的表现。
